% Part: methods
% Chapter: proofs
% Section: proof-by-contradiction

\documentclass[../../../include/open-logic-section]{subfiles}

\begin{document}

\olfileid{mth}{prf}{con}

\olsection{برهان خلف}

در وهلهٔ نخست، برهان خلف الگویی استنتاجی برای اثبات گزاره‌های منفی
است. فرض کنید می‌خواهید نشان دهید گزاره‌ای مانند~$p$ \emph{کاذب} است؛
یعنی می‌خواهید~$\lnot p$ را نشان دهید. امیدبخش‌ترین راهبرد این است که
(الف) فرض کنید $p$ صادق است و (ب) نشان دهید این فرض به چیزی می‌انجامد
که می‌دانید کاذب است. «چیزی که می‌دانیم کاذب است» ممکن است نتیجه‌ای
ناسازگار با خود~$p$ باشد، یعنی آن را نقض کند، یا با فرض دیگری از حکم
کلی مورد بررسی ناسازگار باشد. برای نمونه، اثبات «اگر $q$، آنگاه
$\lnot p$» مستلزم آن است که صدق~$q$ را فرض و~$\lnot p$ را از آن اثبات
کنیم. اگر $\lnot p$ را به برهان خلف اثبات کنید، یعنی افزون بر~$q$،
$p$ را نیز فرض می‌کنید. اگر بتوانید $\lnot q$ را از $p$ اثبات کنید،
نشان داده‌اید که فرض~$p$ به چیزی می‌انجامد که با فرض دیگر شما، یعنی
~$q$، در تناقض است؛ زیرا $q$ و $\lnot q$ نمی‌توانند هر دو صادق باشند.
البته در اثبات تناقض، افزون بر بازکردن تعریف‌ها، باید الگوهای استنتاجی
دیگری را نیز به کار ببرید. مثالی را بررسی کنیم.

\begin{prop}
  اگر $A \subseteq B$ و $B = \emptyset$، آنگاه $A$ هیچ !!{element}ی ندارد.
\end{prop}

\begin{proof}
  فرض کنید $A \subseteq B$ و $B = \emptyset$. می‌خواهیم نشان دهیم
  $A$ هیچ !!{element}ی ندارد.
  \begin{quote}
    چون این یک گزارهٔ شرطی است، مقدم را فرض می‌کنیم و می‌خواهیم تالی
    را اثبات کنیم. تالی چنین است: $A$ هیچ !!{element}ی ندارد. می‌توانیم
    این را اندکی صریح‌تر کنیم: چنین نیست که $x \in A$ای وجود داشته باشد.
  \end{quote}
  $A$ هیچ !!{element}ی ندارد اگر و تنها اگر چنین نباشد که $x$ای وجود
  داشته باشد که $x \in A$.
  \begin{quote}
    پس دریافتیم آنچه واقعاً می‌خواهیم اثبات کنیم گزارهٔ منفی~$\lnot p$
    است، یعنی: چنین نیست که $x \in A$ای وجود داشته باشد. برای استفاده
    از برهان خلف، باید گزارهٔ مثبت متناظر~$p$، یعنی وجود $x \in A$، را
    فرض کنیم و از آن تناقضی به دست آوریم. با نوشتن «به قصد رسیدن به
    تناقض، فرض کنید» یا حتی فقط «خلاف آن را فرض کنید» نشان می‌دهیم که
    برهان خلف به کار می‌بریم و سپس فرض~$p$ را بیان می‌کنیم.
  \end{quote}
  خلاف آن را فرض کنید: $x \in A$ای وجود دارد.
  \begin{quote}
    اکنون این همان فرض تازه‌ای است که برای رسیدن به تناقض به کار
    می‌بریم. دو فرض دیگر هم داریم: $A \subseteq B$ و
    $B = \emptyset$. فرض نخست به ما $x \in B$ را می‌دهد:
  \end{quote}
  چون $A \subseteq B$، داریم $x \in B$.
  \begin{quote}
    اما چون $B = \emptyset$، هر !!{element} از $B$، برای مثال $x$،
    باید !!a{element} از~$\emptyset$ نیز باشد.
  \end{quote}
  چون $B = \emptyset$، داریم $x \in \emptyset$. این تناقض است، زیرا
  بنا بر تعریف، $\emptyset$ هیچ !!{element}ی ندارد.
  \begin{quote}
    همین امر اثبات را کامل می‌کند: از فرض‌هایی که چیده بودیم به چیزی
    که نیاز داشتیم، یعنی تناقض، رسیده‌ایم؛ پس این فرض‌ها نمی‌توانند
    همگی صادق باشند. چون دو فرض نخست، $A \subseteq B$ و
    $B = \emptyset$، محل مناقشه نیستند، باید آخرین فرض افزوده‌شده، یعنی
    وجود $x \in A$، کاذب باشد. بااین‌حال، اگر بخواهیم دقیق و کامل
    باشیم، می‌توانیم این نتیجه را آشکارا بیان کنیم.
  \end{quote}
  بنابراین فرض ما مبنی بر وجود $x \in A$ باید کاذب باشد؛ پس به برهان
  خلف، $A$ هیچ !!{element}ی ندارد.
\end{proof}

هر گزارهٔ مثبت به‌سادگی با گزاره‌ای منفی هم‌ارز است: $p$ اگر و تنها
اگر $\lnot\lnot p$. پس برهان خلف را می‌توان برای اثبات «غیرمستقیم»
گزاره‌های مثبت نیز چنین به کار برد: برای اثبات~$p$، آن را به‌صورت
گزارهٔ منفی $\lnot\lnot p$ بخوانید. اگر بتوانیم از $\lnot p$ تناقضی
به دست آوریم، $\lnot\lnot p$ را به برهان خلف اثبات کرده‌ایم و در نتیجه
~$p$ را به دست آورده‌ایم.

در مثال پیش، هدفمان اثبات گزاره‌ای منفی بود، یعنی اینکه $A$ هیچ
!!{element}ی ندارد؛ بنابراین فرضی که برای برهان خلف ساختیم، یعنی وجود
$x \in A$، گزاره‌ای مثبت بود. این فرض چیزی برای کارکردن به ما داد:
$x \in A$ فرضی‌ای که استدلال دربارهٔ آن را ادامه دادیم تا به
$x \in \emptyset$ رسیدیم.

هنگام اثبات غیرمستقیم یک گزارهٔ مثبت، فرضی که برای برهان خلف می‌سازید
منفی خواهد بود. اما اغلب می‌توانید یک گزارهٔ مثبت را به‌آسانی به
گزاره‌ای منفی، و یک گزارهٔ منفی را به گزاره‌ای مثبت بازنویسی کنید.
اگر در اثبات پیش به جای تالی منفی «$A$ هیچ !!{element}ی ندارد»،
«$A = \emptyset$» را اثبات می‌کردیم، اثبات اساساً همان بود. بنا بر
تعریف $=$، «$A = \emptyset$» یک گزارهٔ کلی است، زیرا به این صورت باز
می‌شود: «هر !!{element} از~$A$، !!{element} از~$\emptyset$ است و
برعکس». اما به‌سادگی می‌توان دید که این گزاره با گزارهٔ منفی «چنین
نیست که $x \in A$ای وجود داشته باشد» هم‌ارز است.

پس گاهی کارکردن با $\lnot p$ به‌عنوان فرض آسان‌تر از اثبات مستقیم~$p$
است. حتی وقتی اثبات مستقیم به همان اندازه ساده یا ساده‌تر است، مانند
مثال‌های بعد، برخی افراد روش غیرمستقیم را ترجیح می‌دهند. اگر نقیض مضاعف
شما را گیج می‌کند، برهان خلف یک گزاره را اثبات تناقضی از
\emph{گزارهٔ مقابل} آن در نظر بگیرید. بنابراین برهان خلف $\lnot p$،
اثبات تناقضی از فرض~$p$ است؛ و برهان خلف~$p$، اثبات تناقضی از
~$\lnot p$ است.

\begin{prop}
$A \subseteq A \cup B$.
\end{prop}

\begin{proof}
  می‌خواهیم نشان دهیم $A \subseteq A \cup B$.
  \begin{quote}
    در نگاه نخست، این گزاره مثبت است: هر $x \in A$ در $A \cup B$ نیز
    هست. نقیض آن چنین است: $x \in A$ای وجود دارد که
    $x \notin A \cup B$. پس می‌توانیم حکم را به‌طور غیرمستقیم اثبات
    کنیم: این گزارهٔ منفی را فرض کنیم و نشان دهیم به تناقض می‌انجامد.
  \end{quote}
  خلاف آن را فرض کنید؛ یعنی $A \nsubseteq A \cup B$.
  \begin{quote}
    برای $A \subseteq A \cup B$ تعریفی داریم: هر $x \in A$ در
    $A \cup B$ نیز هست. برای فهم معنای $A \nsubseteq A \cup B$ باید
    بر تعریف بازشده چند تبدیل منطقی مقدماتی اعمال کنیم: کاذب‌بودن این
    گزاره که هر $x \in A$ در $A \cup B$ نیز هست، اگر و تنها اگر
    \emph{برخی} $x \in A$ها در $A \cup B$ نباشند. اینجا باید بسیار
    دقیق بود: بسیاری از تلاش‌های دانشجویان برای برهان خلف شکست می‌خورد،
    چون نقیض گزاره‌ای مانند «همهٔ $A$ها $B$ هستند» را نادرست تحلیل
    می‌کنند. به بیان دیگر، $A \nsubseteq A \cup B$ اگر و تنها اگر $x$ای
    وجود داشته باشد که $x \in A$ و $x \notin A \cup B$. پس از آن کار
    آسان است.
  \end{quote}
  بنابراین $x \in A$ای وجود دارد که $x \notin A \cup B$. بنا بر تعریف
  $\cup$، $x \in A \cup B$ اگر و تنها اگر $x \in A$ یا $x \in B$.
  چون $x \in A$، داریم $x \in A \cup B$. این با فرض
  $x \notin A \cup B$ تناقض دارد.
\end{proof}

\begin{prob}
\emph{به‌طور غیرمستقیم} اثبات کنید که $A \cap B \subseteq A$.
\end{prob}

\begin{prop}
اگر $A \subseteq B$ و $B \subseteq C$، آنگاه $A \subseteq C$.
\end{prop}

\begin{proof}
  فرض کنید $A \subseteq B$ و $B \subseteq C$. می‌خواهیم نشان دهیم
  $A \subseteq C$.
  \begin{quote}
    به‌طور غیرمستقیم پیش می‌رویم: نقیض آنچه می‌خواهیم اثبات کنیم را
    فرض می‌کنیم.
  \end{quote}
  خلاف آن را فرض کنید؛ یعنی $A \nsubseteq C$.
  \begin{quote}
    مانند پیش استدلال می‌کنیم که $A \nsubseteq C$ اگر و تنها اگر چنین
    نباشد که هر $x \in A$ در $C$ نیز هست؛ یعنی $x \in A$ای وجود داشته
    باشد که $x \notin C$. نگران نباشید؛ با تمرین دیگر بازکردن چنین
    نقیض‌هایی نیازمند فکر زیاد نخواهد بود.
  \end{quote}
  به بیان دیگر، $x$ای وجود دارد که $x \in A$ و $x \notin C$.
  \begin{quote}
    اکنون می‌توانیم از این برای رسیدن به تناقض استفاده کنیم. البته باید
    دو فرض دیگر را نیز به کار ببریم.
  \end{quote}
  چون $A \subseteq B$، داریم $x \in B$. چون $B \subseteq C$، داریم
  $x \in C$. اما این با $x \notin C$ تناقض دارد.
\end{proof}

\begin{prop}
اگر $A \cup B = A \cap B$، آنگاه $A = B$.
\end{prop}

\begin{proof}
  فرض کنید $A \cup B = A \cap B$. می‌خواهیم نشان دهیم $A = B$.
  \begin{quote}
    آغاز کار اکنون معمول است:
  \end{quote}
  به قصد رسیدن به تناقض، فرض کنید $A \neq B$.
  \begin{quote}
    فرض ما برای برهان خلف $A \neq B$ است. چون $A = B$ اگر و تنها اگر
    $A \subseteq B$ و $B \subseteq A$، نتیجه می‌گیریم $A \neq B$ اگر و
    تنها اگر $A \nsubseteq B$ \emph{یا} $B \nsubseteq A$. توجه کنید
    دقت هنگام دست‌کاری نقیض‌ها چقدر مهم است!{} برای به‌دست‌آوردن تناقض
    از این فصل، اثبات به تفکیک حالت‌ها را به کار می‌بریم و نشان می‌دهیم
    در هر حالت تناقضی حاصل می‌شود.
  \end{quote}
  $A \neq B$ اگر و تنها اگر $A \nsubseteq B$ یا $B \nsubseteq A$.
  حالت‌ها را جدا می‌کنیم.
  \begin{quote}
    در حالت نخست $A \nsubseteq B$ را فرض می‌کنیم؛ یعنی برای $x$ای،
    $x \in A$ اما $x \notin B$. $A \cap B$ به‌صورت !!{element}هایی
    تعریف می‌شود که $A$ و $B$ مشترک دارند؛ پس اگر چیزی در یکی از آن‌ها
    نباشد، در اشتراک نیست. $A \cup B$، $A$ همراه با $B$ است؛ پس هرچه در یکی از
    آن دو باشد در اجتماع نیز هست. بنابراین $x \in A \cup B$ اما
    $x \notin A \cap B$، و در نتیجه $A \cap B \neq A \cup B$.
  \end{quote}

  حالت ۱: $A \nsubseteq B$. آنگاه برای $x$ای، $x \in A$ اما
  $x \notin B$. چون $x \notin B$، داریم $x \notin A \cap B$. چون
  $x \in A$، داریم $x \in A \cup B$. پس
  $A \cap B \neq A \cup B$، که با فرض $A \cap B = A \cup B$ تناقض دارد.

  حالت ۲: $B \nsubseteq A$. آنگاه برای $y$ای، $y \in B$ اما
  $y \notin A$. مانند پیش $y \in A \cup B$ ولی
  $y \notin A \cap B$؛ پس باز هم $A \cap B \neq A \cup B$، که با
  $A \cap B = A \cup B$ تناقض دارد.
\end{proof}

\end{document}
